\documentclass[12pt]{article}
\usepackage{fancyhdr}
\pagestyle{fancy}
\fancyhf{}
\fancyfoot[C]{\thepage}
\fancyfoot[L]{\scriptsize CC BY-NC-SA 4.0 \textbullet\ Leo Liang}
\renewcommand{\headrulewidth}{0pt}
\usepackage[margin=1in]{geometry}
\usepackage{amsmath, amssymb}
\usepackage{array, booktabs}
\usepackage{pgffor}
\parindent=0pt
\parskip=1em
\newcolumntype{C}[1]{>{\centering\arraybackslash}p{#1}}
\newcommand{\wl}{\par\vspace{5mm}\noindent\rule{\linewidth}{0.4pt}\par\vspace{1mm}}
\newcommand{\wls}[1]{\foreach \i in {1,...,#1}{\wl}}
% ─────────────────────────────────────────────────────────────────────────────
\begin{document}
\pagestyle{fancy}
\begin{center}
  {\Large\bfseries Homework 6}\\[4pt]
  Domain \& Range: Finding Outputs by Solving Backwards\\[4pt]
  {\small Functions Unit --- Extension}\\[6pt]
  Name:\;\underline{\hspace{6.5cm}}\quad
\end{center}
\hrule\bigskip

\section*{Instructions}
Use a \textbf{pen} and \textbf{show every step}. This is a difficult homework, you must be resourceful. It would be a good idea to draft your answers on a blank sheet of paper before writing your final answer in pen. 

For each function you will find the \textbf{domain} and the
\textbf{range}, and for both you must \emph{justify} your answer using the method below. You may use the internet and AI as a
\textbf{teacher}, but try the problems yourself first. 


% ===========================================================================
\section*{Reading --- Finding the Range by Solving Backwards}
Recall from the functions unit: the \textbf{domain} is the set of allowed inputs, and the \textbf{range} is the
set of outputs the function can actually produce. Finding the domain is about avoiding trouble; finding the range
takes the clever idea you saw in class.

\subsection*{Domain: avoid the two dangers}
For every function in this homework, the domain is all real numbers \emph{except} where the rule breaks. Two
things break a rule:
\begin{itemize}
  \item \textbf{Division by zero:} a denominator can never equal $0$.
  \item \textbf{Even roots of negatives:} whatever sits inside a square root must be $\ge 0$.
\end{itemize}
Everywhere else, the input is allowed.

\subsection*{Range: ask ``which outputs are reachable?''}
Here is the key idea. Pick an \emph{arbitrary} target output $y$ and ask: is there an input $x$ in the domain with
$f(x) = y$? To find out, \textbf{solve the equation $y = f(x)$ for $x$}.
\begin{itemize}
  \item If you can solve it and get a valid $x$ (a real number in the domain), then $y$ \emph{is} in the range.
  \item If certain values of $y$ force an impossible $x$ (dividing by zero, or a square root of a negative), then
        those $y$ are \emph{not} in the range.
\end{itemize}

\textit{This is exactly the ``onto'' question from our functions unit: a value $y$ is in the range precisely when
it has a \textbf{preimage} --- an input that maps to it. Solving $y = f(x)$ for $x$ is how you hunt for that
preimage.}

\subsection*{Worked example}
Find the domain and range of $\;g(x) = \dfrac{6x}{x+2}$.

\textbf{Domain.} The denominator $x+2$ equals $0$ when $x = -2$, so the domain is all real numbers except $-2$.

\textbf{Range.} Set $y = \dfrac{6x}{x+2}$ and solve for $x$:
\begin{align*}
  y(x+2) &= 6x \\
  yx + 2y &= 6x \\
  yx - 6x &= -2y \\
  x(y-6) &= -2y \\
  x &= \frac{-2y}{\,y-6\,}.
\end{align*}
This produces a valid $x$ for every $y$ \emph{except} $y = 6$ (which would divide by zero). So $g$ can output any
real number except $6$: the range is all real numbers except $6$. You can also write this as $(-\infty,6) \cup (6,\infty)$. You can answer in words or use mathematical symbols, either is fine as long as you are clear.

\medskip
Now use this same method on each problem: \textbf{domain first, then solve $y = f(x)$ for $x$}.

% ===========================================================================
\newpage
\section*{Problems}
For each function, find the domain and range with formal justification.

\bigskip
\textbf{1.}\quad $f(x) = 4x - 7$.

\newpage

\bigskip
\textbf{2.}\quad $f(x) = \sqrt{2x + 6}$.

\newpage

\newpage
\textbf{3.}\quad $g(x) = \dfrac{5x}{\,x - 3\,}$.

\newpage

\bigskip
\textbf{4.}\quad $h(x) = x^2 + 2$.

\newpage

\newpage
\textbf{5.}\quad $\bigstar$\quad \textbf{Challenge.}\quad $p(x) = x^2 - 4x + 7$.

\medskip
\textit{Hint: solving $y = x^2 - 4x + 7$ for $x$ means solving a quadratic equation for $x$. Think about when that
quadratic has a real solution --- that is the question that decides the range. (Completing the square is another
route.)}


% ===========================================================================
\newpage
\section*{Notebook Artifact}

\bigskip
\textbf{N1.}\quad If you used the internet or AI, describe how you used it.
\wls{5}

\end{document}